\documentclass[11pt,letterpaper]{exam}

\usepackage{amsmath}
\usepackage{amssymb}
\usepackage{graphicx}
\usepackage{geometry}
\usepackage{longtable}
\geometry{margin=1in}

\begin{document}

\begin{center}
\LARGE \textbf{Midterm Exam - Introduction to Statistics} \\
\large Midterm \\
\normalsize Duration: 75 minutes
\end{center}

\noindent \textbf{Instructions:} Answer all questions. Show your work for full credit. Use the formula sheet provided. You may use a calculator.

\vspace{0.25in}

\begin{questions}

\question[5] What measure of central tendency is most affected by extreme values (outliers)?

\begin{choices}
  \choice Mean
  \choice Median
  \choice Mode
  \choice Range
\end{choices}


\question[5] If a dataset has a standard deviation of 0, what can we conclude?

\begin{choices}
  \choice All values in the dataset are the same
  \choice The dataset has no outliers
  \choice The mean equals the median
  \choice The data is normally distributed
\end{choices}


\question[8] What is the probability of getting at least one head when flipping a fair coin twice?

\begin{choices}
  \choice 0.25
  \choice 0.50
  \choice 0.75
  \choice 1.00
\end{choices}


\question[8] In hypothesis testing, a Type I error occurs when:

\begin{choices}
  \choice We reject a true null hypothesis
  \choice We fail to reject a false null hypothesis
  \choice We accept a true alternative hypothesis
  \choice We reject a false alternative hypothesis
\end{choices}


\question[5] Which of the following $p$-values provides the strongest evidence against the null hypothesis?

\begin{choices}
  \choice $p = 0.001$
  \choice $p = 0.05$
  \choice $p = 0.10$
  \choice $p = 0.50$
\end{choices}


\question[8] If two events $A$ and $B$ are independent, then $P(A \cap B) = $

\begin{choices}
  \choice $P(A) \times P(B)$
  \choice $P(A) + P(B)$
  \choice $P(A) - P(B)$
  \choice $P(A) / P(B)$
\end{choices}


\question[12] Explain the difference between a population parameter and a sample statistic. Provide one example of each.

\vspace{3in}


\question[12] A dataset has a mean of 50 and a standard deviation of 10. Using the empirical rule (68-95-99.7 rule), what percentage of data falls between 40 and 60?

\vspace{3in}


\question[15] Explain what a $p$-value represents in the context of hypothesis testing. Why do we typically use $\alpha = 0.05$ as a significance level?

\vspace{3in}


\question[22] A pharmaceutical company claims their new drug reduces blood pressure by an average of 10 mmHg. Design a hypothesis test to evaluate this claim. Include: (a) null and alternative hypotheses, (b) what Type I and Type II errors would mean in this context, and (c) which error would be more serious and why.

\textit{Hypotheses stated correctly (6 pts), Type I error explained (5 pts), Type II error explained (5 pts), Analysis of error severity (6 pts)}

\vspace{5in}


\newpage

\section*{Formula Sheet}

Sample mean: $\bar{x} = \frac{1}{n}\sum_{i=1}^{n} x_i$

Sample variance: $s^2 = \frac{1}{n-1}\sum_{i=1}^{n}(x_i - \bar{x})^2$

Standard deviation: $s = \sqrt{s^2}$

Z-score: $z = \frac{x - \mu}{\sigma}$

Empirical Rule: 68% within 1 SD, 95% within 2 SD, 99.7% within 3 SD

Probability rules:
- $P(A \cup B) = P(A) + P(B) - P(A \cap B)$
- If independent: $P(A \cap B) = P(A) \times P(B)$

Test statistic: $t = \frac{\bar{x} - \mu_0}{s/\sqrt{n}}$

\end{questions}

\end{document}
